% =============================================================================
% tgp_nbody_lagrangian_eom.tex
% Fragment do \input{} w głównym tomie (np. dodatek D / sekcja n-body).
% Notacja zsynchronizowana z: nbody/pairwise.py, nbody/eom_tgp.py,
% nbody/three_body_force_exact.py, nbody/dynamics_backends.py.
% =============================================================================

\subsection{Equations of motion: $N$ point sources (pairwise TGP + irreducible three-body term)}
\label{subsec:tgp-nbody-lagrangian-eom}

\paragraph{Mass--source identification.}
For point sources in the package \texttt{nbody/} we adopt the same axiom as in the code:
inertial mass equals source strength, $m_i \equiv C_i$ (dimensionless strengths in the units used by the scripts).

\paragraph{Pairwise potential (exact closed form).}
For distinct bodies $i\neq j$ at separation $d_{ij}=|\mathbf{x}_i-\mathbf{x}_j|$, the two-body interaction energy is (cf.\ \texttt{V\_eff} in \texttt{pairwise.py})
\begin{equation}
  V_{2,ij}(d_{ij})
  = -\frac{4\pi\,C_i C_j}{d_{ij}}
  + \frac{8\pi\,\beta\,C_i C_j}{d_{ij}^2}
  - \frac{12\pi\,\gamma\,C_i C_j\,(C_i+C_j)}{d_{ij}^3}.
\end{equation}
The total two-body energy is $V_2=\sum_{i<j}V_{2,ij}(d_{ij})$.

\paragraph{Irreducible three-body term.}
For each unordered triplet $\{i,j,k\}$ the irreducible three-body sector receives
contributions from both $\Phi^3$ and $\Phi^4$, giving
\begin{equation}
  V_{3,ijk}
  = (2\beta - 6\gamma)\,C_i C_j C_k\,
    I(d_{ij},d_{ik},d_{jk}),
  \label{eq:V3-I}
\end{equation}
where $I$ is the triple-overlap integral of single-source Yukawa profiles.  In the Coulomb (unscreened) limit used in closed benchmarks,
\begin{equation}
  I = \frac{8\pi^2}{P},
  \qquad
  P = d_{ij}+d_{ik}+d_{jk}.
\end{equation}
For general screening mass $m_{\rm sp}$ one uses the Feynman two-parameter representation $I_Y(\cdots;m_{\rm sp})$ (see \texttt{three\_body\_force\_exact.py} and \texttt{tgp\_yukawa\_fourier\_feynman\_reduction.tex}).  In the vacuum sector $\beta=\gamma$ of the dimensionless package, the linearized screening mass is $m_{\rm sp}=\sqrt{3\gamma-2\beta}=\sqrt{\gamma}$ (\texttt{tgp\_field.screening\_mass}).

\paragraph{Lagrangian and equations of motion.}
With Cartesian positions $\mathbf{x}_i\in\mathbb{R}^3$ and velocities $\mathbf{v}_i=\dot{\mathbf{x}}_i$,
\begin{equation}
  L = \sum_i \frac{C_i}{2}\,|\mathbf{v}_i|^2 - V(\{\mathbf{x}\}),
  \qquad
  V = V_2 + \sum_{i<j<k} V_{3,ijk}.
\end{equation}
The Euler--Lagrange equations are
\begin{equation}
  C_i\,\ddot{\mathbf{x}}_i = \mathbf{F}_i,
  \qquad
  \mathbf{F}_i = -\nabla_{\mathbf{x}_i} V.
\end{equation}
For a single triplet $(i,j,k)$ contributing \eqref{eq:V3-I}, writing $I=I(d_{ij},d_{ik},d_{jk})$,
\begin{equation}
  \nabla_{\mathbf{x}_i} I
  = \frac{\partial I}{\partial d_{ij}}\,\frac{\mathbf{x}_i-\mathbf{x}_j}{d_{ij}}
  + \frac{\partial I}{\partial d_{ik}}\,\frac{\mathbf{x}_i-\mathbf{x}_k}{d_{ik}},
\end{equation}
and therefore
\begin{equation}
  \mathbf{F}_i^{(ijk)}
  = (6\gamma - 2\beta) C_i C_j C_k\,\nabla_{\mathbf{x}_i} I,
\end{equation}
with analogous expressions for $\mathbf{F}_j^{(ijk)}$ and $\mathbf{F}_k^{(ijk)}$ (pairs involving $j$ and $k$).  Summing over all triplets and adding the pairwise gradient gives the total $\mathbf{F}_i$.

\paragraph{Translational invariance.}
The internal potential depends only on inter-body distances, hence
$\sum_i \mathbf{F}_i=\mathbf{0}$ (total momentum is conserved in isolated $N$-body
subsystems without external fields).  Numerically:
\texttt{dynamics\_backends.total\_force\_from\_accelerations} and
\texttt{examples/ex141\_net\_force\_translational\_invariance.py}.

\paragraph{Rotational invariance and angular momentum.}
If $V$ depends only on Euclidean distances $d_{ij}$, then $V$ is invariant under
a common rigid rotation $\mathbf{x}_i\mapsto R\,\mathbf{x}_i$ for all $i$
($R\in SO(3)$).  Noether's theorem yields conservation of the total angular
momentum (about the origin of the inertial frame)
\begin{equation}
  \mathbf{L}=\sum_i C_i\,\mathbf{x}_i\times\dot{\mathbf{x}}_i,
  \qquad
  \dot{\mathbf{L}}=\sum_i \mathbf{x}_i\times\mathbf{F}_i=\mathbf{0}.
\end{equation}
The Yukawa overlap $I_Y(d_{12},d_{13},d_{23};m)$ is a function of three scalar
edge lengths only, hence does not break $SO(3)$ either.  Numerical check:
\texttt{dynamics\_backends.total\_torque\_about\_origin\_from\_accelerations}
and \texttt{examples/ex143\_net\_torque\_so3\_invariance.py}.
\emph{Remark:} choosing a different reference point $\mathbf{x}_0$ for torques
redefines $\mathbf{L}$ by an overall shift; the statement above is standard for
$\mathbf{x}_i$ measured from an inertial origin with only internal forces.

\paragraph{Conserved momenta along a trajectory.}
For a symplectic time step (e.g.\ leapfrog in \texttt{dynamics\_v2.py}), $\mathbf{P}$
and $\mathbf{L}$ fluctuate at machine-noise level on short horizons when
$(\mathbf{F},\mathbf{V})$ are consistent; see
\texttt{examples/ex144\_conserved\_P\_L\_leapfrog.py}.
Consistency $\mathbf{F}_i=-\nabla_{\mathbf{x}_i}V$ can be probed by central
differences of $V$ (\texttt{forces\_from\_potential\_central\_diff};
\texttt{ex145} (\texttt{coulomb\_3b}), \texttt{ex146} (\texttt{yukawa\_feynman})).

\paragraph{Center of mass and Galilean split.}
Define the total ``inertial mass'' (sum of source strengths)
\begin{equation}
  M_{\rm tot}=\sum_i C_i,
  \qquad
  \mathbf{R}_{\rm cm}=\frac{1}{M_{\rm tot}}\sum_i C_i\,\mathbf{x}_i,
  \qquad
  \mathbf{V}_{\rm cm}=\frac{1}{M_{\rm tot}}\sum_i C_i\,\dot{\mathbf{x}}_i
  =\frac{\mathbf{P}}{M_{\rm tot}}.
\end{equation}
Internal coordinates can be taken as differences $\mathbf{x}_i-\mathbf{x}_j$ or
as $\mathbf{s}_i=\mathbf{x}_i-\mathbf{R}_{\rm cm}$ with the constraint
$\sum_i C_i\,\mathbf{s}_i=\mathbf{0}$.
Because $V$ depends only on distances, $V$ is invariant under a common translation
$\mathbf{x}_i\mapsto \mathbf{x}_i+\mathbf{a}$ for all $i$, hence
$\sum_i\mathbf{F}_i=\mathbf{0}$ and
\begin{equation}
  \ddot{\mathbf{R}}_{\rm cm}
  =\frac{1}{M_{\rm tot}}\sum_i C_i\,\ddot{\mathbf{x}}_i
  =\frac{1}{M_{\rm tot}}\sum_i\mathbf{F}_i=\mathbf{0}.
\end{equation}
Thus the center of mass moves inertially; all nontrivial interaction dynamics
lives in the $(3N-3)$--dimensional quotient by translations (before discussing
rotational reduction).

\paragraph{Hamiltonian form.}
With conjugate momenta $\mathbf{p}_i=\partial L/\partial\dot{\mathbf{x}}_i=C_i\,\dot{\mathbf{x}}_i$,
\begin{equation}
  H=\sum_i\frac{|\mathbf{p}_i|^2}{2C_i}+V(\{\mathbf{x}\}),
  \qquad
  \dot{\mathbf{x}}_i=\frac{\partial H}{\partial\mathbf{p}_i}=\frac{\mathbf{p}_i}{C_i},
  \qquad
  \dot{\mathbf{p}}_i=-\frac{\partial H}{\partial\mathbf{x}_i}=\mathbf{F}_i.
\end{equation}
Symplectic integrators (leapfrog in positions/velocities) are consistent with this
split when implemented as kicks on $\mathbf{v}_i$ and drifts on $\mathbf{x}_i$ with
the same mass identification $m_i=C_i$.
Numerical check of $\ddot{\mathbf{R}}_{\rm cm}\approx\mathbf{0}$:
\texttt{dynamics\_backends.center\_of\_mass\_acceleration} and
\texttt{examples/ex147\_com\_acceleration\_zero.py}.

\paragraph{Split: analytic structure vs.\ coordinate evaluation.}
The \emph{analytic} problem is the closed form above: $V_2$ explicit, $V_3$ through $I$ and its partial derivatives $\partial I/\partial d_{ab}$.
The \emph{numerical} problem is: choose a backend to evaluate $I$ and $\partial I/\partial d_{ab}$ (Feynman quadrature, 3D grid, or Coulomb closure), and separately integrate $\ddot{\mathbf{x}}_i=\mathbf{F}_i/C_i$ in time (symplectic leapfrog or high-order RK as in \texttt{dynamics\_v2.py}).

\paragraph{Lyapunov exponents (optional).}
Tangent-flow Benettin estimates along trajectories are implemented in \texttt{lyapunov.py}; for Hamiltonian flow a symplectic base map (leapfrog + coupled tangent update) is preferred when checking the spectrum sum $\sum\lambda_i\approx 0$ at finite time.
See \texttt{tgp\_lyapunov\_benettin.tex}, scripts \texttt{ex148+} and \texttt{examples/verify\_nbody\_lyapunov\_quick.py}.

\paragraph{Implementation map.}
\begin{itemize}
  \item Pairwise forces: \texttt{analytical\_forces\_tgp\_pairwise} (\texttt{dynamics\_v2.py}).
  \item Coulomb $I=8\pi^2/P$: \texttt{eom\_tgp.py} (\texttt{irreducible\_three\_body\_forces\_from\_I\_derivatives}).
  \item Exact Yukawa $I_Y$: \texttt{three\_body\_force\_exact.py}.
  \item Integrator glue: \texttt{dynamics\_backends.build\_tgp\_integration\_pair}.
  \item Lyapunov / chaos (P1): \texttt{lyapunov.py}, \texttt{tgp\_lyapunov\_benettin.tex}, \texttt{ex148}--\texttt{ex194}, \texttt{verify\_nbody\_lyapunov\_quick.py}.
  \item Short demos: \texttt{ex138}--\texttt{ex140} (integrators); \texttt{ex141}--\texttt{ex143} (identities); \texttt{ex144}--\texttt{ex146} (trajectory / $-\nabla V$); \texttt{ex147} (COM); \texttt{verify\_nbody\_eom\_quick.py}.
\end{itemize}
